% chapters/22_agent_memory.tex
% Part of: AI / LLM / Agents / Hardware Notes

\section{Agent Memory: Episodic, Semantic, Working}

\subsection{In-context (working) memory}

The model's context window is its working memory: it holds the full conversation history,
tool results, and system instructions for the current turn. Everything not in the context
window is, from the model's perspective, forgotten.

Practical implications observed in production agents:
\begin{itemize}
  \item \textbf{Context pressure.} As a session grows, input tokens increase and latency
    / cost rise. Prompt caching mitigates this for stable prefixes (system prompts, skills
    documentation) but not for the growing conversation tail.
  \item \textbf{Shell-snapshots.} Claude Code writes \texttt{shell-snapshots/} entries
    to preserve shell state (environment variables, working directory) between Bash tool
    calls, effectively extending working memory to include filesystem cursor state.
  \item \textbf{Extended thinking.} Thinking blocks give the model a scratchpad within
    a single turn that is not part of the official assistant message — an intra-turn
    working memory separate from the conversational context.
\end{itemize}

\subsection{External episodic memory}

Episodic memory stores a timestamped record of what happened, so the agent can recall
prior sessions or survive context-window resets.

\textbf{JSONL conversation logs} (as implemented in Claude Code) are the clearest
real-world example. Each event in a session is appended as a JSON line to
\texttt{\textasciitilde/.claude/projects/\{project-slug\}/\{session-uuid\}.jsonl}.
Key design choices:

\begin{itemize}
  \item \textbf{Append-only.} New events are always appended; nothing is deleted or
    mutated. This makes the file a reliable audit log and enables trivial crash recovery.

  \item \textbf{Linked-list structure.} The \texttt{parentUuid} field on each record
    points to the preceding event, forming a singly-linked list (or a tree when
    conversations branch via \texttt{isSidechain: true}). Replaying from the head
    reconstructs the full conversation.

  \item \textbf{Turn grouping via \texttt{promptId}.} All events in a single user turn
    (the human message, multiple assistant chunks, and all tool results) share a
    \texttt{promptId}, enabling efficient retrieval of a complete turn without walking
    the full chain.

  \item \textbf{Rich metadata.} Each record stores \texttt{cwd}, \texttt{gitBranch},
    \texttt{version}, \texttt{userType}, and token-usage breakdowns, turning the log
    into a queryable audit trail for debugging and cost analysis.

  \item \textbf{Separation from semantic memory.} A sibling \texttt{memory/} directory
    holds structured facts the agent has been explicitly asked to remember --- distinct
    from the raw episodic log. This mirrors the episodic/semantic distinction in
    cognitive science.
\end{itemize}

\subsection{Semantic memory via vector DBs}

Semantic memory stores structured factual knowledge (facts, rules, definitions) in a
form the agent can query across sessions. Unlike episodic memory, it has no inherent
chronology — it stores \emph{what is true}, not \emph{what happened when}.

\textbf{Five memory types from the CoALA framework} (Cognitive Architectures for
Language Agents, Princeton):

\begin{itemize}
  \item \textbf{Short-term (working):} rolling buffer inside the context window.
    Coherent across a session, forgotten when the window rolls over.
  \item \textbf{Long-term:} persists across sessions via external storage. Crucial for
    personalisation, domain expertise, and continuity. RAG is the primary LTM technique.
  \item \textbf{Episodic:} records of specific past events — actions taken, outcomes
    observed. Enables case-based reasoning: ``last time I saw this error, X fixed it.''
    Implemented by logging event sequences (e.g.\ JSONL session logs).
  \item \textbf{Semantic:} structured factual knowledge. Used in legal/medical diagnosis
    and enterprise knowledge management. Implemented via knowledge graphs,
    vector embeddings, or symbolic AI.
  \item \textbf{Procedural:} skills and learned behaviours enabling automatic task
    execution. Implemented via reinforcement learning or LoRA fine-tuning on task
    demonstrations. Analogous to human muscle memory.
\end{itemize}

\textbf{Vector database as semantic store:}
\begin{itemize}
  \item Facts, documents, and conversation summaries are embedded into dense vectors
    and stored in a vector DB (Chroma, Pinecone, Weaviate, pgvector).
  \item At query time, the agent embeds its current query and retrieves the top-$k$
    semantically similar records via ANN search, injecting them into the context.
  \item \textbf{Personalisation loop:} as the agent accumulates interactions, new
    user preferences, project facts, and domain knowledge are embedded and stored.
    Future sessions retrieve only what is relevant to the current task.
  \item \textbf{Balance constraint:} storing too much leads to slow retrieval and
    noisy context; storing too little loses important facts. Relevance scoring and
    TTL expiry policies help prune stale entries.
\end{itemize}

\subsection{Context engineering}

While ``prompt engineering'' focuses on crafting the right instruction for a single
turn, \textbf{context engineering} is the discipline of designing and managing
\emph{everything} placed in the context window across an agent's full lifecycle: system
prompts, tool schemas, retrieved documents, conversation history, and memory injections.

As agents operate over long horizons and with large tool surfaces, context management
becomes a first-class engineering concern. The key resources from Manus's engineering
team summarise the practical lessons:

\textbf{Context window as shared state.} The context is not just instructions --- it
is the primary medium through which the agent perceives the world, maintains continuity,
and communicates with tools. Every byte in the context competes for the model's
``attention budget''.

\textbf{Core context engineering principles:}
\begin{itemize}
  \item \textbf{Explicit state serialisation.} Rather than relying on the model to
    infer state from conversation history, agents should serialise relevant state
    (current plan, completed steps, open tasks) into structured blocks at the start
    of each turn. This prevents context drift as the window grows.

  \item \textbf{Selective history retention.} Keep full fidelity for recent turns
    (last 2--5) and compress or summarise older turns. Tool call results should be
    trimmed to their essential outputs; raw stdout of thousands of lines is wasted
    context.

  \item \textbf{Chunked retrieval injection.} Retrieved documents and memory fragments
    should be injected near the task description (not buried at the end), with explicit
    relevance labels so the model knows why each fragment was included.

  \item \textbf{Tool schema placement.} Tool descriptions compete with task context
    for tokens. For agents with many tools, hierarchical tool loading (load only the
    relevant tool schema for each step) reduces noise and keeps task-relevant tokens
    more salient.

  \item \textbf{KV cache alignment.} For repeated agentic calls sharing a common
    system prompt and tool suite, structure the context so the stable prefix (system
    prompt + tool definitions) is identical across calls. This enables prompt caching
    to amortise the cost of the static portion across many agent steps.
\end{itemize}

\begin{notebox}
  ``Context engineering'' is the term coined by practitioners at Manus (2025) and
  popularised via the \emph{Context Engineering for AI Agents: Lessons from Building
  Manus} blog post. It reframes the craft of building agents away from individual
  prompts toward systematic management of the entire information flow within the
  context window.
\end{notebox}

\textbf{Context engineering 2.0.} As context windows grow to 1M+ tokens, the problem
shifts from fitting information in to filtering what to prioritise. Research on
``lost-in-the-middle'' (Liu et al.\ 2023) shows that models systematically under-attend
to content in the middle of very long contexts. Mitigation strategies include placing
critical facts near the beginning or end, and using structured XML/Markdown delimiters
to make information boundaries salient.

\subsection{Memory consolidation strategies}

Consolidation moves information from episodic logs into more compact, reusable forms:

\begin{itemize}
  \item \textbf{Explicit summarisation.} At the end of a session (or when context
    pressure builds) the agent can be prompted to write a summary of key decisions and
    facts into the \texttt{memory/} store, discarding raw turn-by-turn detail.

  \item \textbf{Selective retrieval.} Rather than replaying the full JSONL log, an agent
    can embed past turns and retrieve only the $k$ most relevant episodes for the current
    task --- combining episodic storage with a vector-DB retrieval step.

  \item \textbf{Plan persistence.} The \texttt{plans/} subdirectory in
    \texttt{.claude/} acts as a form of procedural memory: multi-step plans survive
    context resets and can be resumed in a new session.

  \item \textbf{Trade-off.} Compact semantic memory loses provenance; raw episodic logs
    are complete but expensive to re-inject. Hybrid approaches maintain a compressed
    ``working summary'' alongside the full log.
\end{itemize}

